\documentclass[11pt,a4paper]{article}
\usepackage[T1]{fontenc}
\usepackage[utf8]{inputenc}
\usepackage{lmodern,microtype,amsmath,amssymb}
\usepackage[margin=21mm]{geometry}
\usepackage{booktabs,longtable,tabularx,array,enumitem,needspace,pdflscape}
\usepackage{xcolor,fancyhdr}
\usepackage[colorlinks=true,linkcolor=blue!45!black,urlcolor=blue!45!black,citecolor=blue!45!black]{hyperref}
\hypersetup{pdftitle={Experimental DIS Rivet analyses in HerwigPolarizedPheno},pdfsubject={Catalogue of experimental measurements and implemented analysis coverage}}
\setlength{\parindent}{0pt}
\setlength{\parskip}{5pt}
\setlength{\emergencystretch}{2em}
\setlength{\headheight}{14pt}
\setlist[description]{font=\normalfont\bfseries,leftmargin=0pt,labelindent=0pt,itemsep=3pt,parsep=0pt}
\pagestyle{fancy}
\fancyhf{}
\fancyhead[L]{\small Experimental DIS Rivet analyses}
\fancyhead[R]{\small HerwigPolarizedPheno}
\fancyfoot[C]{\thepage}
\newcommand{\GeV}{\mathrm{GeV}}
\newcommand{\id}[1]{{\small\texttt{\detokenize{#1}}}}
\newcommand{\summaryname}[1]{\texttt{\detokenize{#1}}}
\newcommand{\summaryfile}[1]{\href{https://github.com/apapaefs/HerwigPolarizedPheno/blob/105045bd757aaa6b217e87811dc5f9a22a8d1cef/analyses/rivet/dis/#1.cc}{\texttt{\detokenize{#1.cc}}}}
\begin{document}

\begin{landscape}
\thispagestyle{empty}
\begingroup
\begin{center}
{\Large\bfseries Experimental DIS Rivet analyses: summary\par}
\end{center}
\vspace{-2pt}
\fontsize{9.3}{11.4}\selectfont
\setlength{\tabcolsep}{3pt}
\renewcommand{\arraystretch}{1.14}
\begin{tabularx}{\linewidth}{@{}l>{\raggedright\arraybackslash}Xlcll@{}}
\toprule
\textbf{Rivet name} & \textbf{Observable} & \textbf{arXiv paper} &
\shortstack{\textbf{Energy [GeV]}\\$E_\ell;\sqrt{s_{\ell N}}$} &
\textbf{Implemented beams} & \textbf{Analysis source file}\\
\midrule
\summaryname{COMPASS_2009_I820721} & Deuteron $A_1^{\pi^\pm,K^\pm}$ & \href{https://arxiv.org/abs/0905.2828}{0905.2828} & 160; 17.4 & $\mu^+p,\mu^+n$ & \summaryfile{COMPASS_2009_I820721}\\
\summaryname{COMPASS_2010_I843494} & Inclusive proton $A_1$ & \href{https://arxiv.org/abs/1001.4654}{1001.4654} & 160; 17.4 & $\mu^+p$ & \summaryfile{COMPASS_2010_I843494}\\
\summaryname{COMPASS_2010_I862410} & Proton $A_1^{\pi^\pm,K^\pm}$ & \href{https://arxiv.org/abs/1007.4061}{1007.4061} & 160; 17.4 & $\mu^+p$ & \summaryfile{COMPASS_2010_I862410}\\
\summaryname{COMPASS_2013_I1236358} & Low-$p_T$ inverse slopes & \href{https://arxiv.org/abs/1305.7317}{1305.7317} & 160; 17.4 & $\mu^+p,\mu^+n$ & \summaryfile{COMPASS_2013_I1236358}\\
\summaryname{COMPASS_2014_I1278730} & Unpolarized cosine amplitudes & \href{https://arxiv.org/abs/1401.6284}{1401.6284} & 160; 17.4 & $\mu^+p,\mu^+n$ & \summaryfile{COMPASS_2014_I1278730}\\
\summaryname{COMPASS_2016_I1357198} & Inclusive proton $A_1$ & \href{https://arxiv.org/abs/1503.08935}{1503.08935} & 200; 19.4 & $\mu^+p$ & \summaryfile{COMPASS_2016_I1357198}\\
\summaryname{COMPASS_2017_I1444985} & $\pi^\pm,h^\pm$ multiplicities & \href{https://arxiv.org/abs/1604.02695}{1604.02695} & 160; 17.4 & $\mu^+p,\mu^+n$ & \summaryfile{COMPASS_2017_I1444985}\\
\summaryname{COMPASS_2017_I1483098} & $K^\pm$ multiplicities & \href{https://arxiv.org/abs/1608.06760}{1608.06760} & 160; 17.4 & $\mu^+p,\mu^+n$ & \summaryfile{COMPASS_2017_I1483098}\\
\summaryname{COMPASS_2017_I1501480} & Inclusive deuteron $A_1$ & \href{https://arxiv.org/abs/1612.00620}{1612.00620} & 160; 17.4 & $\mu^+p,\mu^+n$ & \summaryfile{COMPASS_2017_I1501480}\\
\summaryname{COMPASS_2018_I1624692} & $p_T^2$-dependent $h^\pm$ multiplicities & \href{https://arxiv.org/abs/1709.07374}{1709.07374} & 160; 17.4 & $\mu^+p,\mu^+n$ & \summaryfile{COMPASS_2018_I1624692}\\
\summaryname{COMPASS_2020_I1788430} & High-$z$ $\bar p/p$, $K^-/K^+$ ratios & \href{https://arxiv.org/abs/2003.11791}{2003.11791} & 160; 17.4 & $\mu^+p,\mu^+n$ & \summaryfile{COMPASS_2020_I1788430}\\
\summaryname{COMPASS_2025_I2840545} & Hydrogen $\pi^\pm,K^\pm,h^\pm$ multiplicities & \href{https://arxiv.org/abs/2410.12005}{2410.12005} & 160; 17.4 & $\mu^+p$ & \summaryfile{COMPASS_2025_I2840545}\\
\summaryname{COMPASS_2026_I3096394} & Corrected isoscalar multiplicities & \href{https://arxiv.org/abs/2512.23727}{2512.23727} & 160; 17.4 & $\mu^+p,\mu^+n$ & \summaryfile{COMPASS_2026_I3096394}\\
\summaryname{HERMES_2007_I726689} & Inclusive proton $A_1$ & \href{https://arxiv.org/abs/hep-ex/0609039}{hep-ex/0609039} & 27.6; 7.26 & $e^+p$ & \summaryfile{HERMES_2007_I726689}\\
\summaryname{HERMES_2007_I726689_LEGACY} & Proton $A_\parallel(x,Q^2)$ & \href{https://arxiv.org/abs/hep-ex/0609039}{hep-ex/0609039} & 27.6; 7.26 & $e^+p$ & \summaryfile{HERMES_2007_I726689_LEGACY}\\
\summaryname{HERMES_2013_I1111237} & Unpolarized cosine moments & \href{https://arxiv.org/abs/1204.4161}{1204.4161} & 27.6; 7.26 & $e^+p,e^+n$ & \summaryfile{HERMES_2013_I1111237}\\
\summaryname{HERMES_2013_I1208547} & H/D $\pi^\pm,K^\pm$ multiplicities & \href{https://arxiv.org/abs/1212.5407}{1212.5407} & 27.6; 7.26 & $e^\pm p,e^\pm n$ & \summaryfile{HERMES_2013_I1208547}\\
\summaryname{HERMES_2019_I1698889} & Identified-hadron $A_\parallel$ & \href{https://arxiv.org/abs/1810.07054}{1810.07054} & 27.6; 7.26 & $e^+p,e^+n$ & \summaryfile{HERMES_2019_I1698889}\\
\bottomrule
\end{tabularx}

\vspace{5pt}
All 18 implementations are included; the legacy HERMES entry shares its
experimental paper with the main 2007 entry. Paper identifiers and source
filenames are clickable. All source files are relative to
\href{https://github.com/apapaefs/HerwigPolarizedPheno/tree/105045bd757aaa6b217e87811dc5f9a22a8d1cef/analyses/rivet/dis}{\texttt{analyses/rivet/dis/}}
in the pinned \texttt{HerwigPolarizedPheno} repository.

All entries use a fixed target. The energy column gives the nominal laboratory
lepton energy $E_\ell$ followed by the approximate free-lepton--nucleon collision
energy $\sqrt{s_{\ell N}}=\sqrt{m_\ell^2+m_N^2+2m_NE_\ell}$ for a stationary
proton/neutron. It is distinct from the hadronic invariant mass $W$.
The beam column describes implemented inputs; free $p/n$ yields form the
isoscalar/deuteron predictions. The HERMES 2013 multiplicity source accepts
$e^\pm$, with current cards using $e^+$; the COMPASS 2025 measurement combines
$\mu^\pm$, while its implementation uses $\mu^+$.
\endgroup
\vfill
\begin{center}\thepage\end{center}
\end{landscape}
\begin{center}
{\LARGE\bfseries Experimental DIS analyses\\[3pt] in HerwigPolarizedPheno\par}
\vspace{5pt}
{\large Rivet implementation catalogue\par}
30 September 2026
\end{center}

\section{Scope and interpretation}
This catalogue includes every experimental-measurement-linked Rivet analysis
implemented by a \texttt{.cc} file in
\href{https://github.com/apapaefs/HerwigPolarizedPheno/tree/105045bd757aaa6b217e87811dc5f9a22a8d1cef/analyses/rivet/dis}{\texttt{analyses/rivet/dis/}}.
There are \textbf{18 implementations: 13 COMPASS and five HERMES}, corresponding
to 17 distinct experimental publications. The HERMES 2007 legacy entry is a
second projection of the same measurement. The HERMES 2013 azimuthal entry
corresponds to a published measurement but currently provides a diagnostic
without numerical experimental reference data. All measurements use fixed
targets; no H1 or ZEUS implementation occurs in this directory.

The catalogue is pinned to commit
\href{https://github.com/apapaefs/HerwigPolarizedPheno/commit/105045bd757aaa6b217e87811dc5f9a22a8d1cef}{\texttt{105045bd757a}}.
The 76 directory files were checked against the GitHub snapshot.
Each entry describes the \emph{implemented generator-level selection} and
the subset of published observables covered. Every \texttt{.info} file declares
\texttt{Status: UNVALIDATED}. Reference transcription, generator-level closure
and a full detector/trigger/background/response reproduction are distinct
validation tasks; these entries do not establish the last of these.

\subsection*{Variables and observable construction}
With incoming lepton $k$, outgoing lepton $k'$ and target $P$, define
$q=k-k'$, $Q^2=-q^2$, $x=Q^2/(2P\cdot q)$, $y=P\cdot q/(P\cdot k)$
and $W^2=(P+q)^2$. For a hadron, $z=P\cdot p_h/(P\cdot q)$ and $p_T$
is transverse to the virtual-photon direction. Energies and momenta below
are in GeV; $Q^2$, $W^2$ and $p_T^2$ are in $\GeV^2$.
The COMPASS helpers reconstruct DIS from the highest-energy prompt outgoing
$\mu^+$, excluding muons from hadron decays.

The longitudinal analyses fill helicity-resolved yields and, where required,
inverse-depolarization inputs. Postprocessing combines the physical beam
helicities and constructs $A_\parallel$ or $A_1$, with
$A_\parallel=D(A_1+\eta A_2)$; the inclusive and COMPASS $A_1$ estimators
neglect $\eta A_2$. Multiplicities are hadron yields divided by selected DIS
yields and the specified bin widths. Signed POSNLO and NEGNLO contributions
are combined at normalized-bin level before target sums, ratios or fits.
The unpolarized analyses use zero-polarization samples.

For isoscalar/deuteron predictions, the workflow combines free proton and
neutron samples. This approximation does not reproduce the material target
or nuclear binding, Fermi motion and shadowing. Polarized COMPASS deuteron
longitudinal yields additionally use the 0.925 deuteron D-state factor.
\clearpage

\section{Complete inventory}
\small
\begin{longtable}{@{}p{0.48\textwidth}p{0.41\textwidth}r@{}}
\toprule
Rivet identifier & Implemented experimental observable & Ref.\\
\midrule
\endhead
\id{COMPASS_2009_I820721} & $A_{1,d}^{\pi^\pm,K^\pm}(x)$ & \cite{r820721}\\[4pt]
\id{COMPASS_2010_I843494} & $A_1^p(x)$, 160 GeV & \cite{r843494}\\[4pt]
\id{COMPASS_2010_I862410} & $A_{1,p}^{\pi^\pm,K^\pm}(x)$ & \cite{r862410}\\[4pt]
\id{COMPASS_2013_I1236358} & Low-$p_T$ inverse slopes & \cite{r1236358}\\[4pt]
\id{COMPASS_2014_I1278730} & $\cos\phi_h$, $\cos2\phi_h$ amplitudes & \cite{r1278730}\\[4pt]
\id{COMPASS_2016_I1357198} & $A_1^p(x)$, 200 GeV & \cite{r1357198}\\[4pt]
\id{COMPASS_2017_I1444985} & $\pi^\pm,h^\pm$ multiplicities & \cite{r1444985}\\[4pt]
\id{COMPASS_2017_I1483098} & $K^\pm$ multiplicities & \cite{r1483098}\\[4pt]
\id{COMPASS_2017_I1501480} & $A_1^d(x)$ & \cite{r1501480}\\[4pt]
\id{COMPASS_2018_I1624692} & $h^\pm$ multiplicities versus $p_T^2$ & \cite{r1624692}\\[4pt]
\id{COMPASS_2020_I1788430} & High-$z$ $\bar p/p$, $K^-/K^+$ ratios & \cite{r1788430}\\[4pt]
\id{COMPASS_2025_I2840545} & $\pi^\pm,K^\pm,h^\pm$ on hydrogen & \cite{r2840545}\\[4pt]
\id{COMPASS_2026_I3096394} & Corrected isoscalar multiplicities & \cite{r3096394}\\[4pt]
\id{HERMES_2007_I726689} & Inclusive proton $A_1^p(x)$ & \cite{r726689}\\[4pt]
\id{HERMES_2007_I726689_LEGACY} & Proton $A_\parallel(x,Q^2)$, Table 7 & \cite{r726689}\\[4pt]
\id{HERMES_2013_I1111237} & Unpolarized azimuthal moments & \cite{r1111237}\\[4pt]
\id{HERMES_2013_I1208547} & $\pi^\pm,K^\pm$ multiplicities & \cite{r1208547}\\[4pt]
\id{HERMES_2019_I1698889} & Identified-hadron $A_\parallel$ & \cite{r1698889}\\[4pt]
\bottomrule
\end{longtable}
\normalsize

The 2017 COMPASS pion/hadron and kaon multiplicity releases remain separate
entries for reproducibility. The corrected 2026 addendum supersedes these
older releases. The summaries below identify the stored reference source,
rather than inferring data availability from the presence of a reference
\texttt{.yoda.gz} file: some reference YODA are generated by the workflow.

\subsection*{Common COMPASS multiplicity window}
The 2017 pion/hadron and kaon, 2025 hydrogen, and 2026 corrected isoscalar
implementations apply a species-dependent, whole-$z$-bin restriction to both
the hadron numerator and the DIS denominator:
\begin{equation}
 \frac{\sqrt{(12\,\GeV)^2+m_h^2}}{z_{\min}}
 <\nu=\frac{P\cdot q}{M}<
 \frac{\sqrt{(40\,\GeV)^2+m_h^2}}{z_{\max}}.
 \label{eq:nu}
\end{equation}
Here $m_h$ is the pion mass for pions/unidentified hadrons and the kaon mass
for kaons. This condition is imposed for each published cell. The generated
vertex angle is not used as a replacement for the RICH entrance geometry
corrected through the experiment's acceptance procedure. Fixed-beam cells
outside the implemented kinematic domain remain explicitly masked.
\clearpage
\section{COMPASS analyses}
\Needspace{20\baselineskip}
\subsection{\texorpdfstring{Identified-hadron deuteron spin asymmetries (2009)}{Identified-hadron deuteron spin asymmetries (2009)}}
\id{COMPASS_2009_I820721}\quad\cite{r820721}
\begin{description}
\item[Measurement and setup.] $A_{1,d}^{\pi^\pm,K^\pm}(x)$ from longitudinally polarized 160 GeV $\mu^+$ scattering on deuterium. The implementation combines independent free proton and neutron helicity samples.
\item[Implemented selection.] $Q^2>1$, $0.1<y<0.9$, $0.004<x<0.3$; stable $\pi^\pm,K^\pm$ with $0.2<z<0.85$ and $10<p_h<50$. The helper checks the nominal 160 GeV beam. No explicit $W$ or hadron-angle cut occurs in this source.
\item[Reference coverage.] Four species, each with 10 $x$ bins: 40 asymmetry values from \href{https://www.hepdata.net/record/ins820721?version=1}{HEPData v1, Tables 1--4}. Paper-source differences are retained in a frozen audit.
\item[Scope.] Postprocessing forms $A_1$ from ordinary and inverse-$D$ yields with finite-muon-mass $D$ and R1998, neglecting $\eta A_2$. The plugin does not perform the paper\textquotesingle s quark-helicity extraction or its inclusive measurement.
\end{description}
{\small\href{https://github.com/apapaefs/HerwigPolarizedPheno/blob/105045bd757aaa6b217e87811dc5f9a22a8d1cef/analyses/rivet/dis/COMPASS_2009_I820721.cc}{Analysis source}\quad|\quad\href{https://github.com/apapaefs/HerwigPolarizedPheno/blob/105045bd757aaa6b217e87811dc5f9a22a8d1cef/analyses/rivet/dis/COMPASS_2009_I820721.info}{Metadata}\quad|\quad\href{https://inspirehep.net/literature/820721}{INSPIRE record}}
\Needspace{20\baselineskip}
\subsection{\texorpdfstring{Inclusive proton spin asymmetry at 160 GeV (2010)}{Inclusive proton spin asymmetry at 160 GeV (2010)}}
\id{COMPASS_2010_I843494}\quad\cite{r843494}
\begin{description}
\item[Measurement and setup.] $A_1^p(x)$ from longitudinally polarized 160 GeV $\mu^+$ scattering on the proton in an NH$_3$ target. Generation uses a free proton at rest.
\item[Implemented selection.] $Q^2>1$, $0.1<y<0.9$, $0.004<x<0.7$ and finite positive $D$. There is no explicit $W$ or upper-$Q^2$ veto. A separate nested $Q^2>4$ fill is a validation diagnostic.
\item[Reference coverage.] 15 $x$ bins from \href{https://doi.org/10.17182/hepdata.61588.v2/t1}{HEPData Table 1, v2}. The correct paper is arXiv:1001.4654; this catalogue uses a distinct citation from the 2010 SIDIS entry despite their duplicated metadata BibKey.
\item[Scope.] Physical-helicity postprocessing constructs $A_1$ with finite-muon-mass R1998 depolarization and $\eta A_2=0$. The paper\textquotesingle s $g_1$ determination and Bjorken-sum-rule analysis are outside the implemented observable set.
\end{description}
{\small\href{https://github.com/apapaefs/HerwigPolarizedPheno/blob/105045bd757aaa6b217e87811dc5f9a22a8d1cef/analyses/rivet/dis/COMPASS_2010_I843494.cc}{Analysis source}\quad|\quad\href{https://github.com/apapaefs/HerwigPolarizedPheno/blob/105045bd757aaa6b217e87811dc5f9a22a8d1cef/analyses/rivet/dis/COMPASS_2010_I843494.info}{Metadata}\quad|\quad\href{https://inspirehep.net/literature/843494}{INSPIRE record}}

\clearpage
\Needspace{20\baselineskip}
\subsection{\texorpdfstring{Identified-hadron proton spin asymmetries (2010)}{Identified-hadron proton spin asymmetries (2010)}}
\id{COMPASS_2010_I862410}\quad\cite{r862410}
\begin{description}
\item[Measurement and setup.] $A_{1,p}^{\pi^\pm,K^\pm}(x)$ from longitudinally polarized 160 GeV $\mu^+$ scattering on hydrogen/protons. The current campaign supplies free-proton samples.
\item[Implemented selection.] $Q^2>1$, $0.1<y<0.9$, $0.004<x<0.7$; stable $\pi^\pm,K^\pm$ with $0.2<z<0.85$ and $10<p_h<50$. The shared helper checks the nominal 160 GeV beam. No explicit $W$ or hadron-angle cut is imposed.
\item[Reference coverage.] Four species with 12 $x$ bins each: 48 primary asymmetry values. Numerical values and correlations are reconstructed from checksum-pinned paper source, arXiv:1007.4061.
\item[Scope.] External helicity combination uses event-aggregated ordinary, inverse-$D$ and covariance inputs, finite-muon-mass R1998 depolarization and $\eta A_2=0$. This plugin covers the proton SIDIS asymmetries, rather than all older deuteron data or the paper\textquotesingle s global quark-helicity extraction.
\end{description}
{\small\href{https://github.com/apapaefs/HerwigPolarizedPheno/blob/105045bd757aaa6b217e87811dc5f9a22a8d1cef/analyses/rivet/dis/COMPASS_2010_I862410.cc}{Analysis source}\quad|\quad\href{https://github.com/apapaefs/HerwigPolarizedPheno/blob/105045bd757aaa6b217e87811dc5f9a22a8d1cef/analyses/rivet/dis/COMPASS_2010_I862410.info}{Metadata}\quad|\quad\href{https://inspirehep.net/literature/862410}{INSPIRE record}}
\Needspace{20\baselineskip}
\subsection{\texorpdfstring{Low-transverse-momentum inverse slopes (2013)}{Low-transverse-momentum inverse slopes (2013)}}
\id{COMPASS_2013_I1236358}\quad\cite{r1236358}
\begin{description}
\item[Measurement and setup.] Low-$p_T$ inverse-slope parameter $b$ in $A\exp(-p_T^2/b)$ for $h^\pm$, using spin-averaged 160 GeV $\mu^+$ data on isoscalar $^6$LiD. Generation uses a zero-polarization free-$p/n$ mixture.
\item[Implemented selection.] $Q^2>1$, $W^2>25$, $0.1<y<0.9$; the 23 published $(x,Q^2)$ cells lie within $0.0045\le x<0.12$, $1\le Q^2<10$. Hadrons require $0.2\le z<0.8$, $0.1\le p_T<0.85$, with the pion-mass convention for unidentified hadrons.
\item[Reference coverage.] 368 slope values: 23 cells $\times$ eight $z$ bins $\times$ two charges, transcribed from the pinned paper tables. The \href{https://doi.org/10.1140/epjc/s10052-014-3255-y}{2015 erratum} leaves the slopes unchanged.
\item[Scope.] Generalized least squares fits bin-integrated exponentials to 36 signed $p_T^2$ yields per cell, retaining covariance and negative bins. Variance, covariance and yield-significance requirements are Monte Carlo support criteria. Poor shape is flagged; exact replay of the collaboration\textquotesingle s extraction is unproven. Other spectra and fitted parameters are outside this plugin.
\end{description}
{\small\href{https://github.com/apapaefs/HerwigPolarizedPheno/blob/105045bd757aaa6b217e87811dc5f9a22a8d1cef/analyses/rivet/dis/COMPASS_2013_I1236358.cc}{Analysis source}\quad|\quad\href{https://github.com/apapaefs/HerwigPolarizedPheno/blob/105045bd757aaa6b217e87811dc5f9a22a8d1cef/analyses/rivet/dis/COMPASS_2013_I1236358.info}{Metadata}\quad|\quad\href{https://inspirehep.net/literature/1236358}{INSPIRE record}}

\clearpage
\Needspace{20\baselineskip}
\subsection{\texorpdfstring{Unpolarized azimuthal cosine amplitudes (2014)}{Unpolarized azimuthal cosine amplitudes (2014)}}
\id{COMPASS_2014_I1278730}\quad\cite{r1278730}
\begin{description}
\item[Measurement and setup.] $A_{UU}^{\cos\phi_h}$ and $A_{UU}^{\cos2\phi_h}$ for $h^\pm$ versus $x,z,p_T$ and $(x,z,p_T)$, from target-spin-averaged 160 GeV $\mu^+$ data on $^6$LiD. Generation uses zero-polarization $p/n$ samples.
\item[Implemented selection.] $Q^2>1$, $W^2>25$, $0.003<x<0.13$, $0.2<y<0.9$, virtual-photon angle $\theta_{\gamma^*}<60$ mrad; $0.2\le z<0.85$, $0.1\le p_T<1$. Unidentified hadrons use the pion-mass $z$ convention.
\item[Reference coverage.] 480 cosine-amplitude values from pinned paper Tables 2--12, including overlapping one- and three-dimensional projections.
\item[Scope.] Postprocessing fits bin-integrated constant, cosine and nuisance sine terms to 14 of 16 $\phi_h$ bins, excluding the bins adjacent to zero, and divides by mean depolarization factors with full covariance. The published beam-spin sine asymmetry is omitted; a fitted nuisance is not its prediction. Exact numerical replay of the experimental extraction remains unproven.
\end{description}
{\small\href{https://github.com/apapaefs/HerwigPolarizedPheno/blob/105045bd757aaa6b217e87811dc5f9a22a8d1cef/analyses/rivet/dis/COMPASS_2014_I1278730.cc}{Analysis source}\quad|\quad\href{https://github.com/apapaefs/HerwigPolarizedPheno/blob/105045bd757aaa6b217e87811dc5f9a22a8d1cef/analyses/rivet/dis/COMPASS_2014_I1278730.info}{Metadata}\quad|\quad\href{https://inspirehep.net/literature/1278730}{INSPIRE record}}
\Needspace{20\baselineskip}
\subsection{\texorpdfstring{Inclusive proton spin asymmetry at 200 GeV (2016)}{Inclusive proton spin asymmetry at 200 GeV (2016)}}
\id{COMPASS_2016_I1357198}\quad\cite{r1357198}
\begin{description}
\item[Measurement and setup.] $A_1^p(x)$ from longitudinally polarized 200 GeV $\mu^+$ scattering on an NH$_3$ target, represented by a free proton at rest.
\item[Implemented selection.] $1<Q^2\le190$, $W^2>12$, $0.1<y<0.9$, $0.0025<x<0.7$ and finite positive $D$. A separate $Q^2>4$ fill is a validation diagnostic.
\item[Reference coverage.] 17 $x$ bins from \href{https://doi.org/10.17182/hepdata.72819.v1/t2}{HEPData Table 2, v1}. The reference normalizer documents an upper-edge typo correction while retaining the raw source.
\item[Scope.] The primary comparison is $A_1$, constructed with finite-muon-mass R1998 depolarization and $\eta A_2=0$; $g_1$ columns are provenance only. Fixed 200 GeV generation approximates the accepted experimental momentum interval. With NNPDFpol2.0, the related fitted $g_1$ data prevent interpreting this as independent PDF validation.
\end{description}
{\small\href{https://github.com/apapaefs/HerwigPolarizedPheno/blob/105045bd757aaa6b217e87811dc5f9a22a8d1cef/analyses/rivet/dis/COMPASS_2016_I1357198.cc}{Analysis source}\quad|\quad\href{https://github.com/apapaefs/HerwigPolarizedPheno/blob/105045bd757aaa6b217e87811dc5f9a22a8d1cef/analyses/rivet/dis/COMPASS_2016_I1357198.info}{Metadata}\quad|\quad\href{https://inspirehep.net/literature/1357198}{INSPIRE record}}

\clearpage
\Needspace{20\baselineskip}
\subsection{\texorpdfstring{Pion and unidentified-hadron multiplicities (2017)}{Pion and unidentified-hadron multiplicities (2017)}}
\id{COMPASS_2017_I1444985}\quad\cite{r1444985}
\begin{description}
\item[Measurement and setup.] $dM^{\pi^\pm,h^\pm}/dz$ in sparse $(x,y,z)$ cells for 160 GeV muon DIS on an isoscalar $^6$LiD target. Zero-polarization free $p/n$ yields are summed before the ratio.
\item[Implemented selection.] $Q^2>1$, $W>5$, $0.004<x<0.4$, $0.1<y<0.7$; $0.2\le z\le0.85$, $12<p_h<40$, and Eq.~\eqref{eq:nu} with pion mass in both numerator and denominator.
\item[Reference coverage.] \href{https://www.hepdata.net/record/ins1444985?version=1}{HEPData v1, Tables 1--4}: 311 cells for each of four species, giving 1,244 values. Five cells per species are outside nominal 160 GeV support; 306 remain before statistical masks.
\item[Scope.] The 2026 addendum supersedes this pinned older release. Experimental radiative/diffractive correction columns are retained as provenance; they are not applied to the generator. No generated vertex-angle surrogate for RICH geometry is imposed.
\end{description}
{\small\href{https://github.com/apapaefs/HerwigPolarizedPheno/blob/105045bd757aaa6b217e87811dc5f9a22a8d1cef/analyses/rivet/dis/COMPASS_2017_I1444985.cc}{Analysis source}\quad|\quad\href{https://github.com/apapaefs/HerwigPolarizedPheno/blob/105045bd757aaa6b217e87811dc5f9a22a8d1cef/analyses/rivet/dis/COMPASS_2017_I1444985.info}{Metadata}\quad|\quad\href{https://inspirehep.net/literature/1444985}{INSPIRE record}}
\Needspace{20\baselineskip}
\subsection{\texorpdfstring{Charged-kaon multiplicities (2017)}{Charged-kaon multiplicities (2017)}}
\id{COMPASS_2017_I1483098}\quad\cite{r1483098}
\begin{description}
\item[Measurement and setup.] $dM^{K^\pm}/dz$ in sparse $(x,y,z)$ cells for 160 GeV muon DIS on an isoscalar $^6$LiD target, using zero-polarization free-$p/n$ yields.
\item[Implemented selection.] $Q^2>1$, $W>5$, $0.004<x<0.4$, $0.1<y<0.7$; $0.2\le z\le0.85$, $12<p_h<40$, and Eq.~\eqref{eq:nu} with kaon mass in both numerator and denominator.
\item[Reference coverage.] \href{https://www.hepdata.net/record/ins1483098?version=1}{HEPData v1, Tables 1--2}: 309 cells per charge, giving 618 values. Four cells per charge are outside nominal 160 GeV support; 305 remain before statistical masks.
\item[Scope.] Truth kaon identification and same-event numerator/denominator covariance are used. The corrected 2026 release supersedes this older reference. Unsupported cells are masked, and no RICH vertex-angle surrogate or experimental correction factors are applied to Herwig.
\end{description}
{\small\href{https://github.com/apapaefs/HerwigPolarizedPheno/blob/105045bd757aaa6b217e87811dc5f9a22a8d1cef/analyses/rivet/dis/COMPASS_2017_I1483098.cc}{Analysis source}\quad|\quad\href{https://github.com/apapaefs/HerwigPolarizedPheno/blob/105045bd757aaa6b217e87811dc5f9a22a8d1cef/analyses/rivet/dis/COMPASS_2017_I1483098.info}{Metadata}\quad|\quad\href{https://inspirehep.net/literature/1483098}{INSPIRE record}}

\clearpage
\Needspace{20\baselineskip}
\subsection{\texorpdfstring{Final inclusive deuteron spin asymmetry (2017)}{Final inclusive deuteron spin asymmetry (2017)}}
\id{COMPASS_2017_I1501480}\quad\cite{r1501480}
\begin{description}
\item[Measurement and setup.] $A_1^d(x)$ from longitudinally polarized 160 GeV $\mu^+$ DIS on deuterium. Independent free-$p/n$ samples give $\sigma_{UU}=(p+n)/2$ and $\sigma_{LL}=0.925(p+n)/2$.
\item[Implemented selection.] $1<Q^2\le100$, $W^2>16$, $0.1<y<0.9$, $0.004<x<0.7$ and finite positive $D$. A separate nested $Q^2>4$ fill is diagnostic.
\item[Reference coverage.] 15 $x$ bins from \href{https://doi.org/10.17182/hepdata.78374.v1/t1}{HEPData Table 1, v1}. The experimental $g_1$ columns are retained for provenance rather than used as an implemented $g_1$ prediction.
\item[Scope.] $A_1$ uses finite-muon-mass R1998 depolarization and $\eta A_2=0$. The target approximation omits tensor structure and further nuclear effects. Fixed 160 GeV generation approximates the experimental momentum interval. These related $g_1$ data enter NNPDFpol2.0, so a comparison with that fit is not independent PDF validation.
\end{description}
{\small\href{https://github.com/apapaefs/HerwigPolarizedPheno/blob/105045bd757aaa6b217e87811dc5f9a22a8d1cef/analyses/rivet/dis/COMPASS_2017_I1501480.cc}{Analysis source}\quad|\quad\href{https://github.com/apapaefs/HerwigPolarizedPheno/blob/105045bd757aaa6b217e87811dc5f9a22a8d1cef/analyses/rivet/dis/COMPASS_2017_I1501480.info}{Metadata}\quad|\quad\href{https://inspirehep.net/literature/1501480}{INSPIRE record}}
\Needspace{20\baselineskip}
\subsection{\texorpdfstring{Transverse-momentum-dependent multiplicities (2018)}{Transverse-momentum-dependent multiplicities (2018)}}
\id{COMPASS_2018_I1624692}\quad\cite{r1624692}
\begin{description}
\item[Measurement and setup.] $d^2M^{h^\pm}/(dz\,dp_T^2)$ in sparse $(x,Q^2,z,p_T^2)$ cells for 160 GeV muon DIS on an isoscalar target. Generation uses zero-polarization free $p/n$ components.
\item[Implemented selection.] $Q^2>1$, $W^2>25$, $0.003<x<0.4$, $0.1<y<0.9$; $0.2\le z\le0.8$, $0.02\le p_T^2\le3$, subject to released cell bounds. Stable charged hadrons use the pion-mass convention in $z$.
\item[Reference coverage.] \href{https://www.hepdata.net/record/ins1624692?version=1}{Complete HEPData v1 submission}: 4,664 values, 2,332 per charge from 162 tables. This is the audited released count.
\item[Scope.] Postprocessing sums signed contributions and targets before division by the DIS yield and $\Delta z\,\Delta p_T^2$, retaining covariance. The differential multiplicities are covered; the paper\textquotesingle s subsequent exponential/power-law fit parameters are not. Experimental correction factors are not reapplied to generator events.
\end{description}
{\small\href{https://github.com/apapaefs/HerwigPolarizedPheno/blob/105045bd757aaa6b217e87811dc5f9a22a8d1cef/analyses/rivet/dis/COMPASS_2018_I1624692.cc}{Analysis source}\quad|\quad\href{https://github.com/apapaefs/HerwigPolarizedPheno/blob/105045bd757aaa6b217e87811dc5f9a22a8d1cef/analyses/rivet/dis/COMPASS_2018_I1624692.info}{Metadata}\quad|\quad\href{https://inspirehep.net/literature/1624692}{INSPIRE record}}

\clearpage
\Needspace{20\baselineskip}
\subsection{\texorpdfstring{High-$z$ antiproton/proton and kaon ratios (2020)}{High-z antiproton/proton and kaon ratios (2020)}}
\id{COMPASS_2020_I1788430}\quad\cite{r1788430}
\begin{description}
\item[Measurement and setup.] $R_p=M^{\bar p}/M^p$ and $R_K=M^{K^-}/M^{K^+}$ for 160 GeV $\mu^+$ DIS on isoscalar $^6$LiD. Zero-polarization free-$p/n$ yields are combined before charge ratios.
\item[Implemented selection.] $Q^2>1$, $W^2>25$, $0.01<x<0.4$, $0.1<y<1$, $\theta_h<180$ mrad. Protons/antiprotons: $20\le p_h\le60$, $0.5\le z\le1.1$. Kaons: $x\le0.05$, $40\le p_h\le55$, $0.75\le z\le1.05$. Cell upper edges are admitted within $10^{-12}$ numerical tolerance; global event cuts remain strict.
\item[Reference coverage.] 67 values from pinned paper Tables 1--3: 18 $R_p(x,z)$, 34 low-$x$ $R_p(z,p_h)$ and 15 low-$x$ $R_K(z,p_h)$ values. The inherited kaon selection is separately pinned.
\item[Scope.] Published reconstructed-$z$ bin limits are applied to truth $z$; corrected mean $z$ supplies display coordinates. Same-event charge covariance is retained, and overlapping tables remain separate comparison groups. The frozen provenance records no HEPData submission, so paper tables are the numerical authority.
\end{description}
{\small\href{https://github.com/apapaefs/HerwigPolarizedPheno/blob/105045bd757aaa6b217e87811dc5f9a22a8d1cef/analyses/rivet/dis/COMPASS_2020_I1788430.cc}{Analysis source}\quad|\quad\href{https://github.com/apapaefs/HerwigPolarizedPheno/blob/105045bd757aaa6b217e87811dc5f9a22a8d1cef/analyses/rivet/dis/COMPASS_2020_I1788430.info}{Metadata}\quad|\quad\href{https://inspirehep.net/literature/1788430}{INSPIRE record}}
\Needspace{20\baselineskip}
\subsection{\texorpdfstring{Identified and unidentified multiplicities on hydrogen (2025)}{Identified and unidentified multiplicities on hydrogen (2025)}}
\id{COMPASS_2025_I2840545}\quad\cite{r2840545}
\begin{description}
\item[Measurement and setup.] $dM^{\pi^\pm,K^\pm,h^\pm}/dz$ on liquid hydrogen, measured with 160 GeV muon beams of both electric charges. The implementation uses zero-polarization $\mu^+p$ samples only.
\item[Implemented selection.] $Q^2>1$, $W>5$, $0.004<x<0.4$, $0.1<y<0.7$; $0.2\le z\le0.85$, $12<p_h<40$, and the species-dependent Eq.~\eqref{eq:nu} in numerator and denominator.
\item[Reference coverage.] \href{https://www.hepdata.net/record/ins2840545?version=1}{HEPData v1}: 1,804 values. Each charge has 302 cells for pions and unidentified hadrons and 298 for kaons. All released cells have nominal fixed-energy support.
\item[Scope.] Truth species identification, event-aggregated yields and covariance determine multiplicities. The single-charge simulation is an approximation to the combined-charge experimental sample. Improved experimental radiative corrections remain part of reference provenance; detector/RICH or radiative corrections are not applied to Herwig.
\end{description}
{\small\href{https://github.com/apapaefs/HerwigPolarizedPheno/blob/105045bd757aaa6b217e87811dc5f9a22a8d1cef/analyses/rivet/dis/COMPASS_2025_I2840545.cc}{Analysis source}\quad|\quad\href{https://github.com/apapaefs/HerwigPolarizedPheno/blob/105045bd757aaa6b217e87811dc5f9a22a8d1cef/analyses/rivet/dis/COMPASS_2025_I2840545.info}{Metadata}\quad|\quad\href{https://inspirehep.net/literature/2840545}{INSPIRE record}}

\clearpage
\Needspace{20\baselineskip}
\subsection{\texorpdfstring{Corrected isoscalar multiplicities (2026 addendum)}{Corrected isoscalar multiplicities (2026 addendum)}}
\id{COMPASS_2026_I3096394}\quad\cite{r3096394}
\begin{description}
\item[Measurement and setup.] Corrected $dM^{\pi^\pm,K^\pm,h^\pm}/dz$ for the earlier 160 GeV isoscalar measurements. Zero-polarization $\mu^+p$ and $\mu^+n$ yields form the isoscalar prediction.
\item[Implemented selection.] $Q^2>1$, $W>5$, $0.004<x<0.4$, $0.1<y<0.7$; $0.2\le z\le0.85$, $12<p_h<40$, and Eq.~\eqref{eq:nu} with the appropriate hadron mass in both yields.
\item[Reference coverage.] \href{https://www.hepdata.net/record/ins3096394?version=1}{Corrected HEPData v1 release}: 1,834 values. Each charge has 306 pion/hadron cells and 305 kaon cells. All released cells have nominal fixed-energy support.
\item[Scope.] This addendum supersedes the two 2017 multiplicity references and is the current corrected isoscalar comparison. Improved radiative corrections update the experimental data. Statistical or denominator failures can still mask predictions; nominal support alone does not establish data agreement.
\end{description}
{\small\href{https://github.com/apapaefs/HerwigPolarizedPheno/blob/105045bd757aaa6b217e87811dc5f9a22a8d1cef/analyses/rivet/dis/COMPASS_2026_I3096394.cc}{Analysis source}\quad|\quad\href{https://github.com/apapaefs/HerwigPolarizedPheno/blob/105045bd757aaa6b217e87811dc5f9a22a8d1cef/analyses/rivet/dis/COMPASS_2026_I3096394.info}{Metadata}\quad|\quad\href{https://inspirehep.net/literature/3096394}{INSPIRE record}}

\clearpage
\section{HERMES analyses}
\Needspace{20\baselineskip}
\subsection{\texorpdfstring{Inclusive proton spin asymmetry (2007)}{Inclusive proton spin asymmetry (2007)}}
\id{HERMES_2007_I726689}\quad\cite{r726689}
\begin{description}
\item[Measurement and setup.] $A_1^p(x)$ from longitudinally polarized 27.6 GeV $e^+$ DIS on a proton. Four independent physical-helicity samples provide ordinary and inverse-$D$ cross-section inputs.
\item[Implemented selection.] $1\le Q^2\le20$, $W^2>3.24$, $0.10<y\le0.91$, $0.0212\le x\le0.900$, and scattered-lepton angle $0.04\le\theta_\ell\le0.22$ rad. The highest-energy prompt positron defines DIS kinematics. A $Q^2>4$ selection is diagnostic.
\item[Reference coverage.] 15 proton $A_1$ points from \href{https://doi.org/10.17182/hepdata.11211.v1/t14}{HEPData Table 14}, checked against the official pinned HERMES archive.
\item[Scope.] Postprocessing constructs the $A_1$ proxy with R1990 depolarization and neglects $\eta A_2$. This covers the proton subset of a paper also studying deuteron/neutron $g_1$. The circular lepton-angle envelope does not reproduce the year-dependent rectangular spectrometer response.
\end{description}
{\small\href{https://github.com/apapaefs/HerwigPolarizedPheno/blob/105045bd757aaa6b217e87811dc5f9a22a8d1cef/analyses/rivet/dis/HERMES_2007_I726689.cc}{Analysis source}\quad|\quad\href{https://github.com/apapaefs/HerwigPolarizedPheno/blob/105045bd757aaa6b217e87811dc5f9a22a8d1cef/analyses/rivet/dis/HERMES_2007_I726689.info}{Metadata}\quad|\quad\href{https://inspirehep.net/literature/726689}{INSPIRE record}}
\Needspace{20\baselineskip}
\subsection{\texorpdfstring{Auxiliary low-$Q^2$ proton projection (2007 legacy)}{Auxiliary low-Q2 proton projection (2007 legacy)}}
\id{HERMES_2007_I726689_LEGACY}\quad\cite{r726689}
\begin{description}
\item[Measurement and setup.] Born longitudinal $A_\parallel(x,Q^2)$ from the same 27.6 GeV $e^+p$ measurement as the preceding entry. It is an additional implementation, not an additional experimental publication.
\item[Implemented selection.] $0.18<Q^2<20$, $W^2>3.24$, $0.10<y<0.91$, $0.0041\le x<0.9$, and $0.04<\theta_\ell<0.22$ rad. Physical-helicity yields are combined without depolarization reweighting.
\item[Reference coverage.] 45 measured values in 19 $x$ slices from \href{https://doi.org/10.17182/hepdata.11211.v1/t7}{HEPData Table 7}. The per-$x$ integration boundaries in $Q^2$ remain unverified.
\item[Scope.] Use as an auxiliary projection until authoritative cell boundaries are available. Generator scales below 1 GeV are frozen; the eight points with mean $Q^2<1$ are excluded from quantitative pulls/goodness-of-fit. Exact experimental bin integration and perturbative control in those low-$Q^2$ cells are not established.
\end{description}
{\small\href{https://github.com/apapaefs/HerwigPolarizedPheno/blob/105045bd757aaa6b217e87811dc5f9a22a8d1cef/analyses/rivet/dis/HERMES_2007_I726689_LEGACY.cc}{Analysis source}\quad|\quad\href{https://github.com/apapaefs/HerwigPolarizedPheno/blob/105045bd757aaa6b217e87811dc5f9a22a8d1cef/analyses/rivet/dis/HERMES_2007_I726689_LEGACY.info}{Metadata}\quad|\quad\href{https://inspirehep.net/literature/726689}{INSPIRE record}}

\clearpage
\Needspace{20\baselineskip}
\subsection{\texorpdfstring{Unpolarized azimuthal moments (2013 diagnostic)}{Unpolarized azimuthal moments (2013 diagnostic)}}
\id{HERMES_2013_I1111237}\quad\cite{r1111237}
\begin{description}
\item[Measurement and setup.] Direct $\langle\cos\phi_h\rangle$ and $\langle\cos2\phi_h\rangle$ for $h^\pm,\pi^\pm,K^\pm$ on hydrogen/deuterium. The implementation uses zero-polarization 27.6 GeV $e^+$ with free $p/n$ targets; the publication includes both lepton charges.
\item[Implemented selection.] $Q^2>1$, $W^2>10$, $0.023\le x<0.60$, $0.20\le y<0.85$; $x_F>0.2$, $0.20\le z<1.00$, $0.05\le p_T<1.30$. Hadrons/pions require $1<p_h<15$ and kaons $2<p_h<15$. No explicit lepton- or hadron-angle cut is present.
\item[Reference coverage.] The exact $5\times5\times6\times6$ $(x,y,z,p_T)$ grid contains 900 cells per target/species/charge group. Two moments for 12 groups give 21,600 potential values. \textbf{No numerical experimental values or covariance are vendored} at the pinned revision.
\item[Scope.] This remains a data-free generator diagnostic corresponding to the experimental paper. The repository records an unavailable collaboration endpoint during implementation; it does not supply pseudo-data, experimental pulls or goodness-of-fit.
\end{description}
{\small\href{https://github.com/apapaefs/HerwigPolarizedPheno/blob/105045bd757aaa6b217e87811dc5f9a22a8d1cef/analyses/rivet/dis/HERMES_2013_I1111237.cc}{Analysis source}\quad|\quad\href{https://github.com/apapaefs/HerwigPolarizedPheno/blob/105045bd757aaa6b217e87811dc5f9a22a8d1cef/analyses/rivet/dis/HERMES_2013_I1111237.info}{Metadata}\quad|\quad\href{https://inspirehep.net/literature/1111237}{INSPIRE record}}
\Needspace{20\baselineskip}
\subsection{\texorpdfstring{Identified-hadron multiplicities (2013)}{Identified-hadron multiplicities (2013)}}
\id{HERMES_2013_I1208547}\quad\cite{r1208547}
\begin{description}
\item[Measurement and setup.] Charged-pion/kaon multiplicities on hydrogen and deuterium at 27.6 GeV, in five alternative three-dimensional binnings. The source accepts $e^\pm$; current zero-polarization cards use $e^+$ with free $p/n$ targets.
\item[Implemented selection.] $Q^2>1$, $W^2>10$, $0.10<y<0.85$; matched released cells span $0.023\le x<0.60$, $0.2\le z\le1.1$, $0\le p_T\le1.2$. Stable $\pi^\pm,K^\pm$ require $2<p_h<15$. Only the $(z,Q^2)$ grid imposes the cell upper bound $Q^2=15$; there is no common angle or $x_F$ cut.
\item[Reference coverage.] 6,592 retained cells in 40 target/species/binning datasets from the complete pinned vector-meson-subtracted HERMES archive. There are 104 readable projections with 928 points; \href{https://doi.org/10.17182/hepdata.62097.v1}{64 HEPData tables} provide projection checks.
\item[Scope.] Ratios and density normalizations use event-level numerator/denominator covariance; the released dense statistical covariance is preserved. High-$z$ cells remain annotated. Alternative binnings overlap and require separate comparisons in the absence of cross-binning covariance.
\end{description}
{\small\href{https://github.com/apapaefs/HerwigPolarizedPheno/blob/105045bd757aaa6b217e87811dc5f9a22a8d1cef/analyses/rivet/dis/HERMES_2013_I1208547.cc}{Analysis source}\quad|\quad\href{https://github.com/apapaefs/HerwigPolarizedPheno/blob/105045bd757aaa6b217e87811dc5f9a22a8d1cef/analyses/rivet/dis/HERMES_2013_I1208547.info}{Metadata}\quad|\quad\href{https://inspirehep.net/literature/1208547}{INSPIRE record}}

\clearpage
\Needspace{20\baselineskip}
\subsection{\texorpdfstring{Identified-hadron longitudinal asymmetries (2019)}{Identified-hadron longitudinal asymmetries (2019)}}
\id{HERMES_2019_I1698889}\quad\cite{r1698889}
\begin{description}
\item[Measurement and setup.] Primary $A_\parallel$ for proton $\pi^\pm$ and deuteron $\pi^\pm,K^\pm$ at 27.6 GeV, plus auxiliary charge-difference $A_1$. Independent $e^+p/e^+n$ helicity samples model the publication\textquotesingle s $e^\pm$ hydrogen/deuterium data.
\item[Implemented selection.] $Q^2>1$, $W^2>10$, $y<0.85$, $0.023\le x<0.60$, $0.04<\theta_\ell<0.22$ rad; $x_F>0.1$, $0.1<z<0.8$. Most projections require $z>0.2$; only $(x,z)$ retains $0.1<z<0.2$. Proton pions: $4<p_h<13.8$; deuteron pions/kaons: $2<p_h<15$. Unidentified deuteron hadrons use $0.5<p_h<15$.
\item[Reference coverage.] Pinned APS supplement: 765 primary points (54 $x$, 126 $xz$, 108 $xp_T$, 477 $xzp_T$) plus 36 charge-difference $A_1$ values. The $p_T$ grids extend to 2 GeV; $x$/$xz$ yields have no common upper-$p_T$ veto.
\item[Scope.] No hadron-angle acceptance is imposed. The published deuteron $A_\parallel$ already contains its D-state correction; no additional factor is applied. The 180 published spin-cosine amplitudes are excluded from generator overlays because the implemented optional cosine output is an unpolarized yield moment, not a fitted spin asymmetry.
\end{description}
{\small\href{https://github.com/apapaefs/HerwigPolarizedPheno/blob/105045bd757aaa6b217e87811dc5f9a22a8d1cef/analyses/rivet/dis/HERMES_2019_I1698889.cc}{Analysis source}\quad|\quad\href{https://github.com/apapaefs/HerwigPolarizedPheno/blob/105045bd757aaa6b217e87811dc5f9a22a8d1cef/analyses/rivet/dis/HERMES_2019_I1698889.info}{Metadata}\quad|\quad\href{https://inspirehep.net/literature/1698889}{INSPIRE record}}
\clearpage\begin{thebibliography}{99}\raggedright
\bibitem{r820721} COMPASS Collaboration,
\emph{Flavour Separation of Helicity Distributions from Deep Inelastic Muon-Deuteron Scattering},
Phys. Lett. B 680 (2009) 217--224.
\href{https://doi.org/10.1016/j.physletb.2009.08.065}{doi:10.1016/j.physletb.2009.08.065};\quad\href{https://arxiv.org/abs/0905.2828}{arXiv:0905.2828}.

\bibitem{r843494} COMPASS Collaboration,
\emph{The Spin-dependent Structure Function of the Proton $g_1^p$ and a Test of the Bjorken Sum Rule},
Phys. Lett. B 690 (2010) 466--472.
\href{https://doi.org/10.1016/j.physletb.2010.05.069}{doi:10.1016/j.physletb.2010.05.069};\quad\href{https://arxiv.org/abs/1001.4654}{arXiv:1001.4654}.

\bibitem{r862410} COMPASS Collaboration,
\emph{Quark helicity distributions from longitudinal spin asymmetries in muon-proton and muon-deuteron scattering},
Phys. Lett. B 693 (2010) 227--235.
\href{https://doi.org/10.1016/j.physletb.2010.08.034}{doi:10.1016/j.physletb.2010.08.034};\quad\href{https://arxiv.org/abs/1007.4061}{arXiv:1007.4061}.

\bibitem{r1236358} COMPASS Collaboration,
\emph{Hadron Transverse Momentum Distributions in Muon Deep Inelastic Scattering at 160 GeV/c},
Eur. Phys. J. C 73 (2013) 2531.
\href{https://doi.org/10.1140/epjc/s10052-013-2531-6}{doi:10.1140/epjc/s10052-013-2531-6};\quad\href{https://arxiv.org/abs/1305.7317}{arXiv:1305.7317}.

\bibitem{r1278730} COMPASS Collaboration,
\emph{Measurement of azimuthal hadron asymmetries in semi-inclusive deep inelastic scattering off unpolarised nucleons},
Nucl. Phys. B 886 (2014) 1046--1077.
\href{https://doi.org/10.1016/j.nuclphysb.2014.07.019}{doi:10.1016/j.nuclphysb.2014.07.019};\quad\href{https://arxiv.org/abs/1401.6284}{arXiv:1401.6284}.

\bibitem{r1357198} COMPASS Collaboration,
\emph{The Spin Structure Function $g_1^p$ of the Proton and a Test of the Bjorken Sum Rule},
Phys. Lett. B 753 (2016) 18--28.
\href{https://doi.org/10.1016/j.physletb.2015.11.064}{doi:10.1016/j.physletb.2015.11.064};\quad\href{https://arxiv.org/abs/1503.08935}{arXiv:1503.08935}.

\bibitem{r1444985} COMPASS Collaboration,
\emph{Multiplicities of charged pions and unidentified charged hadrons from deep-inelastic scattering of muons off an isoscalar target},
Phys. Lett. B 764 (2017) 1--10.
\href{https://doi.org/10.1016/j.physletb.2016.09.042}{doi:10.1016/j.physletb.2016.09.042};\quad\href{https://arxiv.org/abs/1604.02695}{arXiv:1604.02695}.

\bibitem{r1483098} COMPASS Collaboration,
\emph{Multiplicities of charged kaons from deep-inelastic muon scattering off an isoscalar target},
Phys. Lett. B 767 (2017) 133--141.
\href{https://doi.org/10.1016/j.physletb.2017.01.053}{doi:10.1016/j.physletb.2017.01.053};\quad\href{https://arxiv.org/abs/1608.06760}{arXiv:1608.06760}.

\bibitem{r1501480} COMPASS Collaboration,
\emph{Final COMPASS results on the deuteron spin-dependent structure function $g_1^d$ and the Bjorken sum rule},
Phys. Lett. B 769 (2017) 34--41.
\href{https://doi.org/10.1016/j.physletb.2017.03.018}{doi:10.1016/j.physletb.2017.03.018};\quad\href{https://arxiv.org/abs/1612.00620}{arXiv:1612.00620}.

\bibitem{r1624692} COMPASS Collaboration,
\emph{Transverse-momentum-dependent Multiplicities of Charged Hadrons in Muon-Deuteron Deep Inelastic Scattering},
Phys. Rev. D 97 (2018) 032006.
\href{https://doi.org/10.1103/PhysRevD.97.032006}{doi:10.1103/PhysRevD.97.032006};\quad\href{https://arxiv.org/abs/1709.07374}{arXiv:1709.07374}.

\bibitem{r1788430} COMPASS Collaboration,
\emph{Antiproton over proton and $K^-$ over $K^+$ multiplicity ratios at high z in DIS},
Phys. Lett. B 807 (2020) 135600.
\href{https://doi.org/10.1016/j.physletb.2020.135600}{doi:10.1016/j.physletb.2020.135600};\quad\href{https://arxiv.org/abs/2003.11791}{arXiv:2003.11791}.

\bibitem{r2840545} COMPASS Collaboration,
\emph{Multiplicities of positive and negative pions, kaons and unidentified hadrons from deep-inelastic scattering of muons off a liquid hydrogen target},
Phys. Rev. D 112 (2025) 012002.
\href{https://doi.org/10.1103/q4rb-bhcg}{doi:10.1103/q4rb-bhcg};\quad\href{https://arxiv.org/abs/2410.12005}{arXiv:2410.12005}.

\bibitem{r3096394} COMPASS Collaboration,
\emph{Addendum to multiplicities of charged pions, kaons and unidentified charged hadrons on an isoscalar target measured by COMPASS Collaboration},
Phys. Lett. B 875 (2026) 140266.
\href{https://doi.org/10.1016/j.physletb.2026.140266}{doi:10.1016/j.physletb.2026.140266};\quad\href{https://arxiv.org/abs/2512.23727}{arXiv:2512.23727}.

\bibitem{r726689} HERMES Collaboration,
\emph{Precise determination of the spin structure function $g_1$ of the proton, deuteron and neutron},
Phys. Rev. D 75 (2007) 012007.
\href{https://doi.org/10.1103/PhysRevD.75.012007}{doi:10.1103/PhysRevD.75.012007};\quad\href{https://arxiv.org/abs/hep-ex/0609039}{arXiv:hep-ex/0609039}.

\bibitem{r1111237} HERMES Collaboration,
\emph{Azimuthal distributions of charged hadrons, pions, and kaons produced in deep-inelastic scattering off unpolarized protons and deuterons},
Phys. Rev. D 87 (2013) 012010.
\href{https://doi.org/10.1103/PhysRevD.87.012010}{doi:10.1103/PhysRevD.87.012010};\quad\href{https://arxiv.org/abs/1204.4161}{arXiv:1204.4161}.

\bibitem{r1208547} HERMES Collaboration,
\emph{Multiplicities of charged pions and kaons from semi-inclusive deep-inelastic scattering by the proton and the deuteron},
Phys. Rev. D 87 (2013) 074029.
\href{https://doi.org/10.1103/PhysRevD.87.074029}{doi:10.1103/PhysRevD.87.074029};\quad\href{https://arxiv.org/abs/1212.5407}{arXiv:1212.5407}.

\bibitem{r1698889} HERMES Collaboration,
\emph{Longitudinal double-spin asymmetries in semi-inclusive deep-inelastic scattering of electrons and positrons by protons and deuterons},
Phys. Rev. D 99 (2019) 112001.
\href{https://doi.org/10.1103/PhysRevD.99.112001}{doi:10.1103/PhysRevD.99.112001};\quad\href{https://arxiv.org/abs/1810.07054}{arXiv:1810.07054}.
\end{thebibliography}
\end{document}
